\documentclass[../main.tex]{subfiles}

\begin{document}

\textit{Chương này trình bày thiết kế hệ thống dựa trên các yêu cầu đã phân tích ở chương 2. Các mục lần lượt trình bày sơ đồ lớp thực thể, thiết kế cơ sở dữ liệu, kiến trúc hướng dịch vụ mô tả bằng các sơ đồ UML và thiết kế giao diện cho ba vai trò.}

Các sơ đồ trong chương dùng ký pháp UML 2.5.1~\cite{uml251} và được vẽ bằng PlantUML~\cite{plantuml} (Chương 1). Thiết kế được đối chiếu với mã nguồn ở phiên bản cuối: schema \texttt{thesis\_core} sau 19 migration Alembic và schema \texttt{thesis\_ai} sau 3 migration.

%=====================================================================
\section{Sơ đồ lớp thực thể}\label{sec:c3-lop}
%=====================================================================

Mô hình miền gồm 22 lớp thực thể. Mỗi lớp ánh xạ sang đúng một bảng của schema \texttt{thesis\_core} (mục~\ref{sec:c3-csdl}). Tên lớp viết tiếng Việt không dấu theo kiểu PascalCase; tên thuộc tính viết theo kiểu snake\_case và trùng với tên cột, để việc đối chiếu sơ đồ lớp, sơ đồ thực thể liên kết và mã nguồn được thuận tiện. Các lớp được chia nhóm theo phân hệ nghiệp vụ của chương 2 (Bảng~\ref{tab:c3-nhom-lop}). Hình~\ref{fig:c3-lop-tong-quat} thể hiện toàn bộ 22 lớp cùng bội số của các liên kết, không kèm thuộc tính; thuộc tính của từng lớp được thể hiện qua các cột tương ứng trong sơ đồ thực thể liên kết ở mục~\ref{subsec:c3-erd}.

{\renewcommand{\arraystretch}{1.2}
\begin{longtable}{|p{3.6cm}|p{6.9cm}|p{4.2cm}|}
\caption{Các nhóm lớp thực thể}\label{tab:c3-nhom-lop}\\
\hline
\textbf{Nhóm} & \textbf{Lớp thực thể} & \textbf{Use case liên quan} \\
\hline
\endfirsthead
\hline
\textbf{Nhóm} & \textbf{Lớp thực thể} & \textbf{Use case liên quan} \\
\hline
\endhead
Người dùng và tổ chức đợt & TaiKhoan, SinhVien, GiangVien, DotDoAn, PhanCongHuongDan, DinhHuong, CauHinhNguong, ThongBao & UC0.1--UC0.5, UC0.8, UC0.9 \\
\hline
Đề tài & DeTai & UC1.1, UC1.2 \\
\hline
Đề cương & DeCuong, TaiLieuDoAn, LuotDuyetDeCuong, LuotDuyetBoMon & UC2.1--UC2.5 \\
\hline
Tiến độ & MocTienDo, BaoCaoTienDo & UC3.1--UC3.4 \\
\hline
Bảo vệ & HoiDong, ThanhVienHoiDong, LichBaoVe, DiemBaoVe & UC4.1--UC4.5 \\
\hline
Kho đồ án và kiểm tra trùng lặp & DoAnKho, KetQuaKiemTraTrungLap, ChiTietTrungLap & UC5.1--UC5.4, UC0.6, UC0.7 \\
\hline
\end{longtable}
}

\begin{figure}[H]
    \centering
    \includegraphics[width=\textwidth]{chuong_3/lop_tong_quat.png}
    \caption{Sơ đồ lớp thực thể tổng quát}
    \label{fig:c3-lop-tong-quat}
\end{figure}

Mô hình có bốn đặc điểm chính:
\begin{itemize}
    \item \textit{Tách tài khoản khỏi hồ sơ.} TaiKhoan giữ thông tin đăng nhập và vai trò; SinhVien và GiangVien giữ hồ sơ (mã số, lớp, học vị). Mỗi tài khoản có tối đa một hồ sơ; tài khoản Bộ môn không có hồ sơ riêng.
    \item \textit{Đề tài là trung tâm.} Đề cương, tài liệu nộp, báo cáo tiến độ, lịch bảo vệ, bản ghi kho và kết quả kiểm tra tên đều gắn với một DeTai. Đề tài do GVHD tạo cho một sinh viên trong \textit{lớp hướng dẫn} của mình, tức tập sinh viên mà Bộ môn gán cho giảng viên đó ở đầu đợt qua PhanCongHuongDan (UC0.8). Mỗi đề tài thuộc một đợt và một định hướng trong danh mục DinhHuong do Bộ môn quản lý (UC0.9); mã định hướng là phần đầu của tên tệp đề cương, nên hệ thống kiểm tra được tên tệp có khớp với định hướng của đề tài hay không.
    \item \textit{Một lớp cho mọi tệp nộp.} TaiLieuDoAn biểu diễn một phiên bản tệp của bất kỳ loại nào: đề cương, bản đề cương đã ký, báo cáo tiến độ, báo cáo cuối, tệp kho và tệp đính kèm đề tài. Mỗi lần nộp tạo một phiên bản mới, phiên bản cũ được giữ nguyên để xem lại lịch sử (UC2.5). Nhờ đó việc lưu trữ, trích xuất văn bản và kiểm tra trùng lặp dùng chung một cách xử lý.
    \item \textit{Lớp liên kết tường minh.} Quan hệ nhiều--nhiều được mô hình hóa thành lớp riêng có thuộc tính: PhanCongHuongDan (sinh viên và giảng viên trong một đợt), ThanhVienHoiDong (giảng viên trong hội đồng, kèm vai trò) và DiemBaoVe (giảng viên chấm một lịch bảo vệ).
\end{itemize}

Quy trình duyệt hai cấp dùng hai lớp lượt duyệt: LuotDuyetDeCuong là kết luận của GVHD cho một phiên bản đề cương (đã duyệt hoặc yêu cầu sửa, UC2.2); LuotDuyetBoMon là kết luận của Lãnh đạo bộ môn (phê duyệt hoặc trả lại, UC2.4), khi phê duyệt thì tham chiếu tới bản PDF đề cương đã ký. Phần kiểm tra trùng lặp gồm KetQuaKiemTraTrungLap, là kết quả một lần kiểm tra cho tên đề tài hoặc cho một phiên bản tài liệu, và ChiTietTrungLap, là một dòng chi tiết gồm đồ án trùng, tỷ lệ, đoạn trong tài liệu nộp, đoạn nguồn và vị trí. Bốn ngưỡng cảnh báo (tên đề tài, tài liệu theo một đồ án, tổng mọi đồ án và từng đoạn) lưu ở CauHinhNguong (UC0.4). DoAnKho là bản ghi kho đồ án, được tạo tự động khi Bộ môn chốt kết quả bảo vệ hoặc nhập từ dữ liệu các khóa cũ.

Các thuộc tính có tập giá trị cố định được khai báo thành kiểu liệt kê. Quan trọng nhất là TrangThaiDeTai (\path{NHAP}, \path{DA_PHAN_CONG}), TrangThaiDeCuong (sáu trạng thái, mục~\ref{subsec:c3-trang-thai}), LoaiTaiLieu (\path{DE_TAI}, \path{DE_CUONG}, \path{DE_CUONG_KY_SO}, \path{BAO_CAO_TIEN_DO}, \path{BAO_CAO_CUOI}, \path{KHO}), LoaiKiemTra (\path{TEN_DE_TAI}, \path{TAI_LIEU}) và TrangThaiKiemTra (\path{DANG_KIEM_TRA}, \path{HOAN_TAT}, \path{LOI}). Một số kiểu giữ lại giá trị cũ của luồng nghiệp vụ trước khi điều chỉnh, như \path{CHO_DUYET}, \path{DA_DUYET}, \path{TU_CHOI} của đề tài, để dữ liệu đã có vẫn hợp lệ; hệ thống không còn tạo ra các giá trị đó.

%=====================================================================
\section{Thiết kế cơ sở dữ liệu}\label{sec:c3-csdl}
%=====================================================================

\subsection{Nguyên tắc thiết kế}\label{subsec:c3-thesis-ai}

Hệ thống dùng MySQL 8 với hai schema trên cùng một máy chủ CSDL: \texttt{thesis\_core} thuộc sở hữu của Web API, chứa toàn bộ dữ liệu nghiệp vụ; \texttt{thesis\_ai} thuộc sở hữu của dịch vụ AI, chứa văn bản đã tách đoạn và vector embedding. Không có khóa ngoại hay câu truy vấn nào đi qua hai schema; dữ liệu chỉ trao đổi qua API nội bộ của dịch vụ AI. Ranh giới này là điều kiện để hai dịch vụ được triển khai và thay đổi độc lập, đúng nguyên tắc sở hữu dữ liệu riêng của SOA (Chương 1). Mọi bảng và cột văn bản dùng \path{utf8mb4} với collation \path{utf8mb4_vietnamese_ci}~\cite{mysqlutf8mb4}, và mọi văn bản được chuẩn hóa NFC~\cite{unicodetr15} tại điểm nhập liệu, nên lưu trữ, so sánh và sắp xếp tiếng Việt có dấu đều đúng. Mọi bảng có khóa chính thay thế \texttt{id} kiểu \texttt{BIGINT} tự tăng; khóa tự nhiên như mã số sinh viên được bảo vệ bằng ràng buộc UNIQUE. Kiểu liệt kê được lưu bằng cột \texttt{VARCHAR} kèm ràng buộc \texttt{CHECK} có tên thay cho kiểu \texttt{ENUM}, nên khi thêm giá trị mới, migration chỉ cần thay ràng buộc \texttt{CHECK}; ngoại lệ duy nhất là cột \path{vai_tro} của bảng \path{tai_khoan}. Mỗi schema có bộ migration Alembic~\cite{alembic} riêng.

Mỗi lớp thực thể ánh xạ sang đúng một bảng có tên là tên lớp viết theo kiểu snake\_case (DeTai thành \path{de_tai}, KetQuaKiemTraTrungLap thành \path{ket_qua_kiem_tra_trung_lap}); thuộc tính ánh xạ sang cột cùng tên; liên kết ánh xạ sang khóa ngoại ở phía bội số nhiều. Các bảng \path{phan_cong_huong_dan}, \path{dinh_huong} và \path{luot_duyet_bo_mon} được bổ sung ở migration 0015 và 0016.

Schema \texttt{thesis\_ai} gồm ba bảng. Bảng \path{chunk} lưu một đơn vị văn bản đã tách: tên đề tài (\path{TOPIC}) và tên đồ án kho (\path{ARCHIVE}) là một đơn vị, tài liệu (\path{DOCUMENT}) gồm nhiều đơn vị mức câu sau khi lọc phần mẫu chung. Cặp (\path{loai_nguon}, \path{ma_nguon}) mới xác định một nguồn, vì \path{ma_nguon} là định danh logic bên \texttt{thesis\_core}, không phải khóa ngoại. Bảng \path{embedding} lưu vector \texttt{float32} của một đơn vị theo một mô hình; tên mô hình gồm tên kho mô hình và mã commit trọng số, nên đổi mô hình không làm lẫn vector cũ và mới. Bảng \path{index_metadata} lưu số vector của mỗi mô hình để phát hiện lệch giữa CSDL và chỉ mục trong bộ nhớ. Gọi \texttt{/ingest} lại cho một nguồn sẽ thay toàn bộ đơn vị cũ của nguồn đó trong một transaction. Chỉ mục FAISS~\cite{faiss2021} và chỉ mục vân tay~\cite{schleimer2003} nằm trong bộ nhớ, được dựng lại từ \path{chunk} và \path{embedding} khi dịch vụ khởi động mà không phải mã hóa lại.

\subsection{Sơ đồ thực thể liên kết và các bảng chính}\label{subsec:c3-erd}\label{subsec:c3-rang-buoc}

Sơ đồ thực thể liên kết (ERD) của \texttt{thesis\_core} được chia thành ba hình theo cùng cách chia nhóm của sơ đồ lớp (Hình~\ref{fig:c3-erd-nguoi-dung-de-tai}, \ref{fig:c3-erd-de-cuong-tien-do} và \ref{fig:c3-erd-bao-ve-kho}). Trong các hình, dấu \texttt{*} trước tên cột chỉ cột \texttt{NOT NULL}; \texttt{<<PK>>}, \texttt{<<FK>>}, \texttt{<<UQ>>}, \texttt{<<CHECK>>} là khóa chính, khóa ngoại, ràng buộc UNIQUE và ràng buộc CHECK; đường nối dùng ký pháp chân chim; bảng vẽ nét đứt thuộc nhóm khác, chỉ thể hiện khóa chính để chỉ ra khóa ngoại trỏ tới.
Bốn bảng của luồng nghiệp vụ chính được mô tả ở Bảng~\ref{tab:c3-bang-de-tai}--\ref{tab:c3-bang-ket-qua}, chỉ gồm các cột chính. Cột ``Ràng buộc'' dùng các ký hiệu PK (khóa chính), FK (khóa ngoại), NN (\texttt{NOT NULL}), UQ (\texttt{UNIQUE}) và CHECK.

\begin{figure}[H]
    \centering
    \includegraphics[width=\textwidth]{chuong_3/erd_nguoi_dung_de_tai.png}
    \caption{ERD nhóm người dùng, đợt đồ án và đề tài}
    \label{fig:c3-erd-nguoi-dung-de-tai}
\end{figure}

{\renewcommand{\arraystretch}{1.15}
\small
\begin{longtable}{|>{\raggedright\arraybackslash}p{4.1cm}|>{\raggedright\arraybackslash}p{2.9cm}|>{\raggedright\arraybackslash}p{1.9cm}|p{5.3cm}|}
\caption{Bảng \texttt{de\_tai}}\label{tab:c3-bang-de-tai}\\
\hline
\textbf{Cột} & \textbf{Kiểu} & \textbf{Ràng buộc} & \textbf{Ý nghĩa} \\
\hline
\endfirsthead
\hline
\textbf{Cột} & \textbf{Kiểu} & \textbf{Ràng buộc} & \textbf{Ý nghĩa} \\
\hline
\endhead
\path{dot_do_an_id} & \texttt{BIGINT} & FK, NN & Đợt đồ án chứa đề tài \\
\hline
\path{giang_vien_id} & \texttt{BIGINT} & FK & GVHD tạo và hướng dẫn đề tài \\
\hline
\path{sinh_vien_id} & \texttt{BIGINT} & FK & Sinh viên được giao đề tài \\
\hline
\path{sinh_vien_id_dang_hoat_dong} & \texttt{BIGINT} & Cột sinh, UQ & Bằng \path{sinh_vien_id} khi đề tài đang hoạt động, ngược lại là \texttt{NULL} \\
\hline
\path{ten_de_tai} & \texttt{VARCHAR(255)} & NN & Tên đề tài \\
\hline
\path{dinh_huong_id} & \texttt{BIGINT} & FK & Định hướng đề tài; rỗng ở đề tài tạo trước khi có danh mục định hướng \\
\hline
\path{trang_thai} & \texttt{VARCHAR(20)} & NN, CHECK & \path{NHAP} hoặc \path{DA_PHAN_CONG}, cùng ba giá trị cũ \\
\hline
\end{longtable}
}

\begin{figure}[H]
    \centering
    \includegraphics[width=\textwidth]{chuong_3/erd_de_cuong_tien_do.png}
    \caption{ERD nhóm đề cương và tiến độ}
    \label{fig:c3-erd-de-cuong-tien-do}
\end{figure}

{\renewcommand{\arraystretch}{1.15}
\small
\begin{longtable}{|>{\raggedright\arraybackslash}p{4.1cm}|>{\raggedright\arraybackslash}p{2.9cm}|>{\raggedright\arraybackslash}p{1.9cm}|p{5.3cm}|}
\caption{Bảng \texttt{de\_cuong}}\label{tab:c3-bang-de-cuong}\\
\hline
\textbf{Cột} & \textbf{Kiểu} & \textbf{Ràng buộc} & \textbf{Ý nghĩa} \\
\hline
\endfirsthead
\hline
\textbf{Cột} & \textbf{Kiểu} & \textbf{Ràng buộc} & \textbf{Ý nghĩa} \\
\hline
\endhead
\path{de_tai_id} & \texttt{BIGINT} & FK, NN, UQ & Đề tài sở hữu đề cương; mỗi đề tài có tối đa một đề cương \\
\hline
\path{trang_thai} & \texttt{VARCHAR(20)} & NN, CHECK & Một trong sáu trạng thái đề cương \\
\hline
\path{nop_bo_mon_luc} & \texttt{DATETIME} & & Lần GVHD nộp đề cương lên Bộ môn gần nhất \\
\hline
\path{ket_qua_bo_mon} & \texttt{VARCHAR(20)} & & Kết quả gần nhất của Bộ môn (\path{DA_DUYET} hoặc \path{TRA_LAI}); được xóa khi sinh viên nộp phiên bản mới \\
\hline
\end{longtable}
}

{\renewcommand{\arraystretch}{1.15}
\small
\begin{longtable}{|>{\raggedright\arraybackslash}p{4.1cm}|>{\raggedright\arraybackslash}p{2.9cm}|>{\raggedright\arraybackslash}p{1.9cm}|p{5.3cm}|}
\caption{Bảng \texttt{tai\_lieu\_do\_an}}\label{tab:c3-bang-tai-lieu}\\
\hline
\textbf{Cột} & \textbf{Kiểu} & \textbf{Ràng buộc} & \textbf{Ý nghĩa} \\
\hline
\endfirsthead
\hline
\textbf{Cột} & \textbf{Kiểu} & \textbf{Ràng buộc} & \textbf{Ý nghĩa} \\
\hline
\endhead
\path{de_tai_id} & \texttt{BIGINT} & FK, NN & Đề tài của tài liệu \\
\hline
\path{de_cuong_id} & \texttt{BIGINT} & FK & Đề cương, với tài liệu thuộc quy trình đề cương \\
\hline
\path{loai} & \texttt{VARCHAR(20)} & NN, CHECK & Loại tài liệu, một trong sáu loại \\
\hline
\path{so_phien_ban} & \texttt{INT} & NN & Số phiên bản trong cùng loại của một đề tài \\
\hline
\path{ten_luu_tru} & \texttt{VARCHAR(255)} & NN, UQ & Tên tệp do hệ thống sinh trong kho tệp \\
\hline
\path{dinh_dang} & \texttt{VARCHAR(10)} & NN, CHECK & \texttt{DOCX} hoặc \texttt{PDF}, xác định theo nội dung nhị phân của tệp \\
\hline
\path{noi_dung_trich_xuat} & \texttt{LONGTEXT} & & Văn bản trích xuất đã chuẩn hóa NFC, đầu vào của kiểm tra trùng lặp \\
\hline
\path{da_index_luc} & \texttt{DATETIME} & & Thời điểm tài liệu vào chỉ mục của dịch vụ AI; \texttt{NULL} là chưa vào \\
\hline
\end{longtable}
}

\begin{figure}[H]
    \centering
    \includegraphics[width=0.95\textwidth]{chuong_3/erd_bao_ve_kho_trung_lap.png}
    \caption{ERD nhóm bảo vệ, kho đồ án và kiểm tra trùng lặp}
    \label{fig:c3-erd-bao-ve-kho}
\end{figure}

{\renewcommand{\arraystretch}{1.15}
\small
\begin{longtable}{|>{\raggedright\arraybackslash}p{4.1cm}|>{\raggedright\arraybackslash}p{2.9cm}|>{\raggedright\arraybackslash}p{1.9cm}|p{5.3cm}|}
\caption{Bảng \texttt{ket\_qua\_kiem\_tra\_trung\_lap}}\label{tab:c3-bang-ket-qua}\\
\hline
\textbf{Cột} & \textbf{Kiểu} & \textbf{Ràng buộc} & \textbf{Ý nghĩa} \\
\hline
\endfirsthead
\hline
\textbf{Cột} & \textbf{Kiểu} & \textbf{Ràng buộc} & \textbf{Ý nghĩa} \\
\hline
\endhead
\path{de_tai_id} & \texttt{BIGINT} & FK & Đề tài liên quan \\
\hline
\path{tai_lieu_do_an_id} & \texttt{BIGINT} & FK, UQ & Phiên bản tài liệu được kiểm tra; rỗng với kiểm tra tên và tệp tải lên \\
\hline
\path{tai_khoan_id} & \texttt{BIGINT} & FK & Người tải tệp lên kiểm tra đạo văn; rỗng với tài liệu đã nộp \\
\hline
\path{loai_kiem_tra} & \texttt{VARCHAR(20)} & NN, CHECK & \path{TEN_DE_TAI} hoặc \path{TAI_LIEU} \\
\hline
\path{trang_thai} & \texttt{VARCHAR(20)} & NN, CHECK & \path{DANG_KIEM_TRA}, \path{HOAN_TAT} hoặc \path{LOI} \\
\hline
\path{ty_le_tong} & \texttt{DECIMAL(5,2)} & & Tỷ lệ ký tự bị gắn cờ trùng trên toàn văn bản (\%); giá trị được so với ngưỡng tổng \\
\hline
\path{ty_le_nguon_cao_nhat} & \texttt{DECIMAL(5,2)} & & Tỷ lệ trùng với một nguồn cao nhất; giá trị được so với ngưỡng tài liệu \\
\hline
\path{nguong_ap_dung} & \texttt{DECIMAL(5,2)} & NN & Ngưỡng tại thời điểm kiểm tra \\
\hline
\path{nguong_tong_ap_dung} & \texttt{DECIMAL(5,2)} & & Ngưỡng tổng mọi nguồn tại thời điểm kiểm tra tài liệu \\
\hline
\end{longtable}
}

Ngưỡng áp dụng được lưu cùng kết quả, nên khi Bộ môn đổi ngưỡng, kết quả cũ vẫn hiển thị đúng nhãn vượt hay dưới ngưỡng tại thời điểm kiểm tra. Mỗi dòng \path{chi_tiet_trung_lap} ghi nhãn đoạn (trùng nguyên văn, trùng sửa nhẹ hoặc trùng ngữ nghĩa) và xác định nguồn bằng cặp (\path{loai_nguon}, \path{nguon_id}), không có khóa ngoại vì nguồn có thể thuộc ba bảng khác nhau.

\textbf{Ràng buộc toàn vẹn và xử lý đồng thời.} Các quy tắc nghiệp vụ quan trọng được biểu diễn bằng ràng buộc của CSDL thay vì chỉ kiểm tra ở tầng ứng dụng, vì kiểu kiểm tra rồi mới ghi có thể để lọt cả hai yêu cầu đến gần như đồng thời. Quy tắc mỗi sinh viên có tối đa một đề tài đang hoạt động dùng cột sinh \path{sinh_vien_id_dang_hoat_dong} kèm chỉ mục UNIQUE: MySQL cho phép nhiều giá trị \texttt{NULL} trong chỉ mục UNIQUE, nên ràng buộc chỉ chặn khi cùng một sinh viên xuất hiện ở hai đề tài đang hoạt động, và Web API trả mã 409 cho yêu cầu ghi thứ hai. Ràng buộc UNIQUE trên \path{tai_lieu_do_an_id} của bảng kết quả bảo đảm mỗi phiên bản tài liệu có tối đa một kết quả kiểm tra, nên kết quả về muộn của phiên bản cũ không ghi đè kết quả của phiên bản mới; các ràng buộc UNIQUE khác thể hiện trên ERD. Các thao tác đổi trạng thái được viết thành câu \texttt{UPDATE} có điều kiện trên trạng thái hiện tại rồi kiểm tra số dòng bị ảnh hưởng; thao tác cần đọc nhiều bảng trước khi ghi, như nộp đề cương hay lưu phân công, khóa dòng bằng \texttt{SELECT \ldots{} FOR UPDATE} trong cùng transaction. Riêng quy tắc chỉ chốt điểm bảo vệ khi mọi thành viên hội đồng đã chấm được kiểm tra trong transaction ở tầng ứng dụng.

%=====================================================================
\section{Kiến trúc hướng dịch vụ (SOA)}\label{sec:c3-kien-truc}
%=====================================================================

Hệ thống gồm ba thành phần triển khai độc lập. \textit{Client Web} là ứng dụng React một trang, build thành tệp tĩnh và được Apache phục vụ. \textit{Web API} là ứng dụng FastAPI sở hữu toàn bộ nghiệp vụ và schema \texttt{thesis\_core}; đây là thành phần duy nhất client gọi tới. \textit{Dịch vụ AI} là ứng dụng FastAPI riêng, sở hữu schema \texttt{thesis\_ai}, chịu trách nhiệm kiểm tra trùng lặp tên đề tài và tài liệu. Kiến trúc dựa trên các nguyên tắc và quyết định thiết kế chính sau:
\begin{itemize}
    \item \textit{Một cửa vào công khai.} Trình duyệt chỉ giao tiếp với Apache qua HTTPS. Web API và dịch vụ AI không mở cổng ra ngoài.
    \item \textit{Tách dịch vụ AI thành tiến trình và schema riêng.} Mô hình embedding nặng, nhu cầu tài nguyên khác Web API, và lỗi của phần AI không làm dừng nghiệp vụ. Đánh đổi là thêm một điểm lỗi mạng và dữ liệu phải được đồng bộ qua API.
    \item \textit{Hợp đồng tường minh.} Web API cung cấp REST API có tài liệu OpenAPI sinh tự động; dịch vụ AI cung cấp API nội bộ với các model Pydantic~\cite{pydantic} kiểm tra dữ liệu vào ra.
    \item \textit{Đẩy dữ liệu khi ghi} (\textit{push-on-write}). Sau khi giao đề tài, sau khi kiểm tra một tài liệu hoàn tất và khi thêm bản ghi kho, Web API gọi \texttt{/ingest} trong tác vụ nền; dịch vụ AI không đọc ngược \texttt{thesis\_core}. Chỉ mục vì thế có thể trễ, nên cột \path{da_index_luc} (với tài liệu) và \path{trang_thai_index} (với bản ghi kho) cho biết nguồn chưa vào chỉ mục để Web API thử lại ở các lần sau.
    \item \textit{Kiểm tra tên đồng bộ, kiểm tra tài liệu bất đồng bộ.} Tên ngắn, kết quả cần hiển thị ngay khi giao đề tài. Tài liệu dài, thời gian kiểm tra có thể tới hàng chục giây, nên không giữ kết nối HTTP mà để client hỏi lại trạng thái.
    \item \textit{Lỗi của dịch vụ AI không chặn nghiệp vụ.} Khi dịch vụ AI lỗi hoặc quá thời gian chờ, việc giao đề tài và nộp tài liệu vẫn thành công; kết quả kiểm tra được ghi là lỗi và có thể kiểm tra lại.
\end{itemize}

\subsection{Sơ đồ thành phần và sơ đồ triển khai}

Hình~\ref{fig:c3-thanh-phan} là sơ đồ thành phần của hệ thống. Web API được chia thành các module theo phân hệ nghiệp vụ: \textit{xác thực và phân quyền} phát hành, kiểm tra JWT~\cite{rfc7519} và kiểm tra vai trò của từng endpoint; \textit{quản trị} quản lý tài khoản, đợt đồ án, định hướng, phân công GVHD và cấu hình ngưỡng; năm module \textit{Đề tài, Đề cương, Tiến độ, Bảo vệ, Kho đồ án} hiện thực nghiệp vụ của năm phân hệ ở chương 2; \textit{thông báo và tổng quan}; \textit{lưu trữ tệp} kiểm tra định dạng theo nội dung nhị phân, lưu tệp ngoài thư mục web, trích xuất văn bản từ DOCX/PDF và chuẩn hóa NFC. Module \textit{AI client} là lớp duy nhất gọi dịch vụ AI: lớp này gắn khóa API, đặt thời gian chờ theo từng loại lời gọi và chuyển mọi lỗi kết nối, hết thời gian chờ hay sai hợp đồng thành một kiểu lỗi chung để module nghiệp vụ xử lý. Dịch vụ AI gồm bốn thành phần nối tiếp: API nội bộ; khâu tiền xử lý (chuẩn hóa NFC, lọc phần mẫu chung như mục lục và tài liệu tham khảo, tách câu, tách từ tiếng Việt bằng underthesea~\cite{underthesea}); mô hình embedding \path{vietnamese-bi-encoder}~\cite{bkaibiencoder}; chỉ mục FAISS cho so khớp ngữ nghĩa và chỉ mục vân tay (winnowing) cho so khớp nguyên văn. Chi tiết các kỹ thuật này đã trình bày ở Chương 1.

\begin{figure}[H]
    \centering
    \includegraphics[width=\textwidth]{chuong_3/so_do_thanh_phan.png}
    \caption{Sơ đồ thành phần của hệ thống}
    \label{fig:c3-thanh-phan}
\end{figure}

Hình~\ref{fig:c3-trien-khai} là sơ đồ triển khai ở mức thiết kế. Máy chủ ứng dụng chia thành hai vùng. \textit{Vùng công khai} chỉ có máy chủ Web Apache; Apache kết thúc TLS (phiên bản 1.2 trở lên) ở cổng 443, chuyển cổng 80 sang 443, phục vụ bản build tĩnh của client và chuyển tiếp các yêu cầu \path{/api/} tới Web API bằng \texttt{mod\_proxy}~\cite{apachemodproxy}. \textit{Vùng nội bộ} gồm Web API, dịch vụ AI, MySQL và kho tệp tải lên, không công bố cổng ra ngoài. Không có đường dẫn nào của Apache trỏ tới dịch vụ AI, nên dịch vụ AI chỉ nhận được yêu cầu từ Web API.

\begin{figure}[H]
    \centering
    \includegraphics[width=0.7\textwidth]{chuong_3/so_do_trien_khai.png}
    \caption{Sơ đồ triển khai mức thiết kế}
    \label{fig:c3-trien-khai}
\end{figure}

Client gọi Web API bằng REST/JSON với JWT trong header \texttt{Authorization}; Web API gọi dịch vụ AI bằng REST/JSON nội bộ với khóa trong header \texttt{X-API-Key}. Mỗi dịch vụ truy cập schema của mình bằng tài khoản CSDL riêng, Web API qua SQLAlchemy~\cite{sqlalchemy}; tài khoản phục vụ của dịch vụ AI không có quyền DDL. Kho tệp tải lên nằm ngoài thư mục gốc của Apache: người dùng chỉ tải được tệp qua endpoint của Web API, sau khi Web API kiểm tra quyền xem tài liệu đó. Thiết kế này được hiện thực bằng Docker Compose~\cite{dockercompose}: mỗi thành phần là một container, vùng nội bộ là một mạng Docker riêng, và chỉ container Apache công bố cổng 80 và 443 (chi tiết ở chương 4).

\subsection{Hợp đồng API và sơ đồ tuần tự kiểm tra trùng lặp}

Dịch vụ AI cung cấp năm endpoint (Bảng~\ref{tab:c3-api-ai}). Bốn endpoint nghiệp vụ yêu cầu khóa \texttt{X-API-Key} dài tối thiểu 32 ký tự và được so sánh bằng hàm so sánh thời gian hằng; endpoint \texttt{/health} chỉ trả trạng thái để kiểm tra sống, không cần khóa. Trường \path{loai_tru} của \texttt{/check-document} chứa định danh các phiên bản khác của cùng đề tài: dịch vụ AI không biết khái niệm đề tài, nên Web API gửi tường minh danh sách này để một phiên bản mới không bị coi là trùng với phiên bản trước của chính nó.

{\renewcommand{\arraystretch}{1.15}
\small
\begin{longtable}{|>{\raggedright\arraybackslash}p{2.6cm}|>{\raggedright\arraybackslash}p{3.9cm}|>{\raggedright\arraybackslash}p{4.1cm}|>{\raggedright\arraybackslash}p{3.6cm}|}
\caption{Các endpoint của dịch vụ AI}\label{tab:c3-api-ai}\\
\hline
\textbf{Endpoint} & \textbf{Đầu vào} & \textbf{Đầu ra} & \textbf{Bên gọi, thời điểm} \\
\hline
\endfirsthead
\hline
\textbf{Endpoint} & \textbf{Đầu vào} & \textbf{Đầu ra} & \textbf{Bên gọi, thời điểm} \\
\hline
\endhead
\texttt{POST} \path{/check-title} & \path{text}, \path{top_k}, \path{nguong} & Danh sách \path{matches} gồm \path{ref_id}, \path{loai} (\path{TOPIC} hoặc \path{ARCHIVE}), \path{ty_le} & Web API, đồng bộ khi GVHD nhập hoặc giao đề tài; thời gian chờ 4 giây \\
\hline
\texttt{POST} \path{/check-document} & \path{ref_id}, \path{text}, \path{top_k}, \path{nguong_doan}, \path{nguong_nguyen_van}, \path{loai_tru} & \path{ty_le_tong}, \path{do_an_trung}, \path{chi_tiet} (đoạn nguồn, đoạn trùng, vị trí, nhãn \path{NGUYEN_VAN}, \path{SUA_NHE} hoặc \path{NGU_NGHIA}), \path{so_doan}, \path{so_doan_loai_tru}, \path{ty_le_loai_tru} & Tác vụ nền sau khi nộp tài liệu; thời gian chờ tăng theo độ dài văn bản, tối đa 240 giây \\
\hline
\texttt{POST} \path{/ingest} & \path{id}, \path{type} (\path{TOPIC}, \path{ARCHIVE} hoặc \path{DOCUMENT}), \path{text}, \path{nhom} (mã đề tài, chỉ với \path{DOCUMENT}) & \path{chunk_count} & Tác vụ nền sau khi giao đề tài, sau khi kiểm tra tài liệu hoàn tất, khi thêm bản ghi kho \\
\hline
\texttt{POST} \path{/reindex} & Danh sách mục ingest (tùy chọn) & \path{indexed_count} & Nạp kho hàng loạt; mã hóa bổ sung vector còn thiếu sau khi đổi mô hình \\
\hline
\texttt{GET} \path{/health} & Không có & Trạng thái, tên mô hình, kích thước chỉ mục, số đơn vị thiếu vector & Kiểm tra sống khi khởi động và giám sát \\
\hline
\end{longtable}
}

\textbf{Kiểm tra tên đề tài khi giao đề tài} (Hình~\ref{fig:c3-tuan-tu-ten}). Trong lúc GVHD nhập tên trên MH-06, client gọi \texttt{POST} \path{/api/topics/check-title} để hiển thị cảnh báo ngay; kết quả này không được lưu. Khi GVHD bấm giao đề tài, client gửi \texttt{POST} \path{/api/topics}. Web API kiểm tra vai trò, đợt đang mở và sinh viên thuộc lớp hướng dẫn, chuẩn hóa NFC tên đề tài, đọc ngưỡng tên đề tài rồi gọi \texttt{/check-title}; nếu dịch vụ AI không phản hồi trong 4 giây hoặc trả lỗi, kết quả được ghi là \path{LOI} và việc giao vẫn tiếp tục. Web API ghi đề tài ở trạng thái \path{DA_PHAN_CONG}, kết quả kiểm tra, các dòng chi tiết và thông báo cho sinh viên trong một transaction rồi trả mã 201. Sau đó, tác vụ nền gọi \texttt{/ingest} để đưa tên đề tài mới vào chỉ mục.

\begin{figure}[H]
    \centering
    \includegraphics[width=0.9\textwidth]{chuong_3/tuan_tu_kiem_tra_ten.png}
    \caption{Sơ đồ tuần tự kiểm tra trùng lặp tên đề tài khi giao đề tài}
    \label{fig:c3-tuan-tu-ten}
\end{figure}

\textbf{Kiểm tra tài liệu khi nộp đề cương hoặc báo cáo cuối} (Hình~\ref{fig:c3-tuan-tu-tai-lieu}). Web API kiểm tra định dạng theo nội dung nhị phân và dung lượng tối đa 20~MB, lưu tệp, trích xuất văn bản, rồi ghi phiên bản tài liệu cùng một kết quả kiểm tra ở trạng thái \path{DANG_KIEM_TRA} trong cùng transaction và trả mã 202 ngay. Tác vụ nền gọi \texttt{/check-document} với ngưỡng từng đoạn và danh sách loại trừ. Nếu thành công, kết quả chuyển sang \path{HOAN_TAT}, các dòng chi tiết được ghi và tài liệu được gửi tới \texttt{/ingest} để trở thành nguồn so sánh cho các lần kiểm tra sau; nếu lỗi hoặc hết thời gian chờ, kết quả chuyển sang \path{LOI}. Client hỏi lại trạng thái bằng \texttt{GET} \texttt{/api/documents/\{id\}/check} cho tới khi kết quả không còn ở \path{DANG_KIEM_TRA}. Thiết kế giữ bất biến: bản ghi không ở mãi trạng thái đang kiểm tra. Nếu tiến trình Web API khởi động lại giữa chừng, lần đọc kết quả kế tiếp sẽ chuyển bản ghi đã quá 300 giây sang \path{LOI}; từ \path{LOI}, GVHD hoặc Bộ môn có thể yêu cầu kiểm tra lại.

\begin{figure}[H]
    \centering
    \includegraphics[width=\textwidth]{chuong_3/tuan_tu_kiem_tra_tai_lieu.png}
    \caption{Sơ đồ tuần tự kiểm tra trùng lặp tài liệu khi nộp}
    \label{fig:c3-tuan-tu-tai-lieu}
\end{figure}

\subsection{Sơ đồ trạng thái}\label{subsec:c3-trang-thai}

\textbf{Đề tài} (Hình~\ref{fig:c3-trang-thai-de-tai}). GVHD có thể lưu nháp hoặc giao ngay. Bản nháp chỉ GVHD thấy; khi GVHD gửi, đề tài chuyển sang đã phân công và sinh viên nhận thông báo. Đề tài không còn bước duyệt riêng, vì sinh viên đã thống nhất đề tài với GVHD trước khi GVHD giao trên hệ thống.

\begin{figure}[H]
    \centering
    \includegraphics[width=0.7\textwidth]{chuong_3/trang_thai_de_tai.png}
    \caption{Sơ đồ trạng thái đề tài}
    \label{fig:c3-trang-thai-de-tai}
\end{figure}

\textbf{Đề cương} (Hình~\ref{fig:c3-trang-thai-de-cuong}). Đề cương đi qua hai cấp duyệt: GVHD duyệt nội dung, sau đó Lãnh đạo bộ môn phê duyệt bằng bản đề cương đã ký để nộp về khoa. Sinh viên nộp phiên bản đầu tiên (UC2.1) đưa đề cương từ \path{CHUA_NOP} sang \path{CHO_DUYET}, kích hoạt kiểm tra trùng lặp và thông báo cho GVHD. GVHD duyệt phiên bản mới nhất (\path{GVHD_DA_DUYET}) hoặc yêu cầu sửa kèm nhận xét (\path{YEU_CAU_SUA}, UC2.2), và có thể trả lại đề cương đã duyệt. GVHD nộp lên Bộ môn (\path{CHO_BO_MON_DUYET}, UC2.3) khi tên tệp đúng mẫu và mã định hướng khớp đề tài. Bộ môn phê duyệt bằng cách tải lên bản PDF đề cương đã ký (\path{DA_DUYET}), hoặc trả lại về \path{YEU_CAU_SUA} kèm nhận xét bắt buộc (UC2.4). Từ \path{YEU_CAU_SUA}, sinh viên nộp phiên bản mới để quay lại \path{CHO_DUYET}; phiên bản cũ được giữ. \path{DA_DUYET} là trạng thái cuối; báo cáo tiến độ và bảo vệ chỉ mở cho đề tài có đề cương ở trạng thái này. Hệ thống không đọc hay xác minh chữ ký trong bản PDF tải lên; việc ký thuộc trách nhiệm của Lãnh đạo bộ môn.

\begin{figure}[H]
    \centering
    \includegraphics[width=\textwidth]{chuong_3/trang_thai_de_cuong.png}
    \caption{Sơ đồ trạng thái đề cương}
    \label{fig:c3-trang-thai-de-cuong}
\end{figure}

%=====================================================================
\section{Thiết kế giao diện}\label{sec:c3-giao-dien}
%=====================================================================

Giao diện được thiết kế trước khi thi công dưới dạng prototype HTML/CSS tĩnh (Chương 1). Mỗi màn hình là một trang HTML dùng tên class theo quy ước component của Ant Design~\cite{antd}, nên khi thi công bằng React, mỗi vùng của prototype ứng với một component có sẵn. Ảnh trong mục này được chụp tự động từ prototype bằng trình duyệt Chrome chạy không giao diện, ở độ rộng 1440 điểm ảnh; dữ liệu trong các màn hình đều là dữ liệu hư cấu. Giao diện theo phong cách quản trị của Ant Design và được thiết kế cho máy tính, vì người dùng chính là giảng viên và Bộ môn: menu dọc nền tối bên trái, thanh tiêu đề có đường dẫn điều hướng (breadcrumb), chuông thông báo và tên người dùng; nội dung chính là bảng và biểu mẫu đặt trong thẻ (Card). Mỗi trạng thái nghiệp vụ hiển thị bằng Tag có màu kèm nhãn chữ, nên vẫn đọc được khi in trắng đen. Các trạng thái khác của cùng màn hình (đang kiểm tra, lỗi, hộp thoại) được minh họa trong khối riêng bên dưới màn hình chính.

Ba vai trò dùng chung bố cục, chỉ khác menu. Menu của sinh viên gồm Tổng quan, Đề tài của tôi, Đề cương, Báo cáo tiến độ, Bảo vệ và Kho đồ án. Menu của giảng viên gồm Tổng quan, Sinh viên của tôi, Đề tài hướng dẫn, Giao đề tài, Duyệt đề cương, Nhận xét tiến độ, Hội đồng bảo vệ và Kho đồ án. Menu của Bộ môn gồm Tổng quan toàn khoa, Quản lý đợt đồ án, Phân công GVHD, Đề tài toàn khoa, Duyệt đề cương, Tiến độ toàn khoa, Hội đồng và lịch bảo vệ, Kho đồ án và Cấu hình ngưỡng trùng lặp. Phía máy chủ, Web API vẫn kiểm tra vai trò ở từng endpoint, nên việc ẩn menu chỉ là lớp hỗ trợ sử dụng, không phải cơ chế bảo mật.

Bộ thiết kế gồm 16 màn hình chính (Bảng~\ref{tab:c3-man-hinh}), phủ đủ năm phân hệ và nhóm chức năng chung. Các màn hình quản lý danh mục đơn giản như tài khoản, đợt đồ án, định hướng và mốc tiến độ dùng lại mẫu bảng và biểu mẫu của các màn hình này, nên không thiết kế riêng.

{\renewcommand{\arraystretch}{1.15}
\small
\begin{longtable}{|p{1.6cm}|p{6.5cm}|p{3.5cm}|p{2.6cm}|}
\caption{Danh mục màn hình thiết kế}\label{tab:c3-man-hinh}\\
\hline
\textbf{Mã} & \textbf{Tên màn hình} & \textbf{Vai trò} & \textbf{Use case} \\
\hline
\endfirsthead
\hline
\textbf{Mã} & \textbf{Tên màn hình} & \textbf{Vai trò} & \textbf{Use case} \\
\hline
\endhead
MH-01 & Đăng nhập & Cả ba vai trò & UC0.1 \\
\hline
MH-02 & Tổng quan của sinh viên & Sinh viên & UC0.5, UC4.5 \\
\hline
MH-03 & Tổng quan của giảng viên & Giảng viên & UC1.1, UC0.8 \\
\hline
MH-04 & Tổng quan toàn khoa, kèm tiến độ & Bộ môn & UC3.1, UC3.4 \\
\hline
MH-05 & Danh sách đề tài & Cả ba vai trò & UC1.2 \\
\hline
MH-06 & Giao đề tài, có cảnh báo trùng tên & Giảng viên & UC1.1, UC0.6 \\
\hline
MH-07 & Chi tiết đề tài & Cả ba vai trò & UC1.2 \\
\hline
MH-08 & Nộp đề cương và báo cáo cuối & Sinh viên & UC2.1, UC2.5, UC4.1 \\
\hline
MH-09 & Duyệt đề cương hai cấp & Giảng viên, Bộ môn & UC2.2--UC2.4 \\
\hline
MH-10 & Báo cáo tiến độ và nhận xét & Sinh viên, Giảng viên & UC3.2, UC3.3 \\
\hline
MH-11 & Lập hội đồng và xếp lịch & Bộ môn & UC4.2, UC4.4 \\
\hline
MH-12 & Chấm điểm bảo vệ & Giảng viên & UC4.3 \\
\hline
MH-13 & Kho đồ án & Cả ba vai trò & UC5.1, UC5.3, UC5.4 \\
\hline
MH-14 & Chi tiết kết quả kiểm tra trùng lặp & Giảng viên, Bộ môn & UC0.7 \\
\hline
MH-15 & Phân công GVHD & Bộ môn & UC0.8 \\
\hline
MH-16 & Sinh viên của tôi & Giảng viên & UC0.8, UC1.1 \\
\hline
\end{longtable}
}

Phần tiếp theo trình bày các màn hình gắn với luồng nghiệp vụ chính: trang tổng quan, giao đề tài, duyệt đề cương và kết quả kiểm tra trùng lặp.

\textbf{MH-02 -- Tổng quan của sinh viên} (Hình~\ref{fig:c3-mh-02}). Bốn thẻ cho sinh viên thấy ngay tình trạng của mình: đề tài và GVHD, đề cương (phiên bản, trạng thái duyệt, tỷ lệ trùng lặp gần nhất), tiến độ theo các mốc, lịch và kết quả bảo vệ. Ảnh minh họa mở sẵn danh sách thông báo chưa đọc của chuông thông báo.

\begin{figure}[H]
    \centering
    \includegraphics[width=\textwidth]{chuong_3/mh_02.png}
    \caption{MH-02 -- Tổng quan của sinh viên}
    \label{fig:c3-mh-02}
\end{figure}

\textbf{MH-06 -- Giao đề tài} (Hình~\ref{fig:c3-mh-06}). GVHD nhập tên, chọn sinh viên trong lớp hướng dẫn, chọn định hướng (bắt buộc) và nhập mô tả. Ngay khi nhập tên, khối cảnh báo liệt kê các đề tài tương tự trong kho kèm khóa, GVHD và tỷ lệ tương đồng. Cảnh báo không chặn việc giao: GVHD có thể sửa tên hoặc vẫn giao. GVHD cũng có thể lưu nháp để giao sau.

\begin{figure}[H]
    \centering
    \includegraphics[width=\textwidth]{chuong_3/mh_06.png}
    \caption{MH-06 -- Giao đề tài có cảnh báo trùng tên}
    \label{fig:c3-mh-06}
\end{figure}

\textbf{MH-09 -- Duyệt đề cương} (Hình~\ref{fig:c3-mh-09}). Màn hình phục vụ cả hai cấp duyệt. Phần trên dành cho GVHD: nội dung tóm tắt của phiên bản mới nhất, bảng kiểm tra trước khi nộp lên khoa (tên tệp đúng mẫu, mã định hướng khớp đề tài, trích xuất được nội dung), ô nhận xét và hai nút yêu cầu sửa hoặc duyệt. Bên phải là kết quả kiểm tra trùng lặp với tỷ lệ tổng, ngưỡng áp dụng và các đồ án trùng nhiều nhất. Khối giữa minh họa trạng thái khi kết quả chưa có: GVHD vẫn duyệt được. Khối cuối minh họa hai trạng thái sau: GVHD nộp đề cương đã duyệt lên Bộ môn, và Lãnh đạo bộ môn tải lên bản PDF đề cương đã ký để phê duyệt hoặc trả lại kèm nhận xét.

\begin{figure}[H]
    \centering
    \includegraphics[width=\textwidth]{chuong_3/mh_09.png}
    \caption{MH-09 -- Duyệt đề cương hai cấp}
    \label{fig:c3-mh-09}
\end{figure}

\textbf{MH-14 -- Chi tiết kết quả kiểm tra trùng lặp} (Hình~\ref{fig:c3-mh-14}). Màn hình trình bày kết quả theo kiểu báo cáo của các công cụ kiểm tra đạo văn. Thanh tóm tắt gồm tỷ lệ tổng, tỷ lệ trùng với một đồ án cao nhất, hai ngưỡng (theo một đồ án và tổng mọi đồ án) và nhãn vượt hoặc dưới ngưỡng; khi tài liệu chỉ vượt ngưỡng tổng, nhãn ghi rõ là trùng tổng với nhiều đồ án. Cột trái liệt kê các đồ án trùng và đóng vai trò bộ lọc. Hai cột bên phải đặt đoạn trong tài liệu nộp cạnh đoạn nguồn. Mỗi đoạn có nhãn chữ ghi mức độ (cao, trung bình, thấp) và tỷ lệ, không chỉ dựa vào màu.

\begin{figure}[H]
    \centering
    \includegraphics[width=\textwidth]{chuong_3/mh_14.png}
    \caption{MH-14 -- Chi tiết kết quả kiểm tra trùng lặp}
    \label{fig:c3-mh-14}
\end{figure}

%=====================================================================
\section{Kết luận chương}
%=====================================================================

Chương 3 đã trình bày thiết kế hệ thống ở bốn khía cạnh. Mô hình miền gồm 22 lớp thực thể, chia theo năm phân hệ và nhóm chức năng chung, trong đó đề tài là trung tâm và mọi tệp nộp dùng chung một lớp tài liệu có phiên bản. Cơ sở dữ liệu gồm hai schema tách biệt: \texttt{thesis\_core} với 22 bảng do Web API sở hữu và \texttt{thesis\_ai} với ba bảng do dịch vụ AI sở hữu. Các quy tắc nghiệp vụ quan trọng như mỗi sinh viên có một đề tài đang hoạt động, mỗi phiên bản tài liệu có một kết quả kiểm tra được bảo đảm bằng ràng buộc của CSDL. Kiến trúc hướng dịch vụ gồm client Web, Web API và dịch vụ AI, với một cửa vào công khai qua Apache, dữ liệu được đẩy sang dịch vụ AI khi ghi, kiểm tra tên đồng bộ và kiểm tra tài liệu bất đồng bộ. Các sơ đồ trạng thái mô tả vòng đời đề tài và quy trình duyệt đề cương hai cấp. Mười sáu màn hình thiết kế theo quy ước Ant Design phủ đủ các use case chính của ba vai trò. Chương 4 trình bày việc thi công hệ thống theo thiết kế này.

\end{document}
